% Source: https://artofproblemsolving.com/community/c3h3590758_midpoints_and_foot_of_the_height
% Stable identity: aops-topic-3590758-q01
Similarity gives $CD=\dfrac{1}{2}CA$, so $D$ is the midpoint of the hypotenuse $AC$ of right triangle $APC$. Consequently $DP=DC=DA$. Also $\angle PCD=\angle ACB=20^\circ$, because $P$ lies on segment $BC$. The isosceles triangle $DPC$ has base angles $20^\circ$, so $\angle PDC=180^\circ-40^\circ=\boxed{140^\circ}$.
\begin{asy}
import geometry;
size(240);
pair C=(0,0),B=(6,0),A=5*dir(20),D=A/2,E=B/2,P=(A.x,0);draw(A--B--C--cycle);draw(D--E);draw(A--P^^D--P,dashed);label("$A$",A,N);label("$B$",B,SE);label("$C$",C,SW);label("$D$",D,NW);label("$E$",E,S);label("$P$",P,S);
\end{asy}
